\section*{Chapter \ref{ch:rs:order-book-features} --- Order-Book Features}

\begin{solution}{exo:rs:order-book-features:1}
$I = (700 - 300)/1\,000 = 0.4$; the mid is 99.985; the weighted mid is $(99.99 \times 700 + 99.98 \times 300)/1\,000 = 99.987$, moved towards
the ask because the ask queue is the smaller.
\end{solution}

\begin{solution}{exo:rs:order-book-features:2}
The bid rose, so $+q^b_n = +200$ and the old bid queue is not subtracted; the ask did not change, so $-400 + 400 = 0$: $e_n = +200$.
\end{solution}

\begin{solution}{exo:rs:order-book-features:3}
$11\,000/2\,750 = 4$ average queue sizes; $4 \times 0.47 = 1.88$ ticks.
\end{solution}

\begin{solution}{exo:rs:order-book-features:4}
With $\lambda = \mu$ a queue of $n$ lots takes a time to empty that grows without bound in probability as $n$ grows, while the one-lot queue
empties in a time that does not depend on $n$; so the ask almost surely survives the bid, and P(up) tends to zero. The model's extreme
buckets are therefore near 0 and 1, whereas the data stay between 0.25 and 0.74: something other than the queues' own random walks moves
the price.
\end{solution}

\begin{solution}{exo:rs:order-book-features:5}
At one second most intervals have no mid change, and a change is a single tick set off by one event, of which OFI sees only part; over
longer intervals the changes are sums of many events, their discreteness averages out, and the linear relation with the summed flow shows.
\end{solution}

\begin{solution}{exo:rs:order-book-features:6}
$(a q^b + b q^a)/(q^b + q^a) = \frac{a + b}{2} + \frac{a - b}{2}\cdot\frac{q^b - q^a}{q^b + q^a} = m + \frac s2 I$.
\end{solution}

\begin{solution}{exo:rs:order-book-features:7}
0.277 for the five-level \omterm{def:rs:order-book-features:depth}{depth imbalance} against 0.253 for the touch imbalance: the deeper levels add information here because the
simulator's liquidity providers also cancel stale quotes behind the touch. On real data the deeper levels are also where orders are
cheapest to fake.
\end{solution}

\begin{solution}{exo:rs:order-book-features:8}
The features must be computed at the time the firm could have them, the receive time, not the exchange time; and a consolidated book built
on exchange timestamps assumes all venues' messages arrive at once, which they do not. The accuracy is also uninformative without the
base rate: at extreme imbalances the naive direction is already right three times in four. Recompute with receive-time books and report the
accuracy by bucket against the base rate.
\end{solution}

\begin{solution}{pb:rs:order-book-features:1}
\textbf{1.} 0.245 and 0.736. \textbf{2.} 16\,564 and 14\,538. \textbf{3.} 3.4 against 55.4 lots, and 49.2 against 3.1. \textbf{4.} Near zero,
slightly below (the bucket centred on $-0.1$ gives 0.500). \textbf{5.} 2.03 lots a second arrive and 2.13 leave. \textbf{6.} 0.017 and 0.979.
\textbf{7.} It treats the queues as independent random walks, while the efficient price drives which side is cancelled and traded against,
and new liquidity keeps arriving at the smaller queue. \textbf{8.} The efficient price: informed orders and stale-quote cancellations that
depend on where it is. \textbf{9.} 10\%, 53\% and 79\%; 0.47 ticks per average best-queue size. \textbf{10.} 2\,747 shares (27.5 lots).
\textbf{11.} A mean squared error 10\% below the mid's. \textbf{12.} The weighted mid: 12\% and 3\% below the mid. \textbf{13.} 0.30 and
0.25. \textbf{14.} Because a model researched on one set of numbers and traded on another is a different model; any discrepancy is a
silent bug. \textbf{15.} 2\,766 two-sided rows of 2\,769 messages. \textbf{16.} OFI and the rates need every event in order; the imbalances,
the weighted mid and the \omterm{def:rs:order-book-features:micro}{microprice} can be computed from snapshots. \textbf{17.} Depth beyond the touch, and the touch imbalance itself:
orders placed without intent to trade. \textbf{18.} \emph{Named result:} across tenths of \omterm{def:rs:order-book-features:imbalance}{queue imbalance} the probability of an up move runs
from 0.245 to 0.736 on the simulated tape, where the race of two birth--death queues predicts 0.017 to 0.979; OFI explains 53\% of ten-second
mid changes. \textbf{19.} The efficient price: arrivals and cancellations that depend on it (a queue-reactive model with a hidden state).
\textbf{20.} Where supply and demand stand right now, and, with its flow, where the price is about to go.
\end{solution}

\begin{solution}{iq:rs:order-book-features:1}
Most likely up: the small ask queue is much closer to emptying, and liquidity providers on the ask may be the ones withdrawing. But not with
the certainty a naive queue race suggests; on the simulated tape the most bid-heavy tenth gives about 0.74, not 0.98.
\iqlookfor{direction plus calibration: imbalance is informative but noisy.}
\end{solution}

\begin{solution}{iq:rs:order-book-features:2}
Net quantity added at the bid, minus removed from it, minus added at the ask, plus removed from it, counting whole queues when prices change.
A net order flow of size $x$ consumes or builds queues of typical size $D$; the price moves by about one tick per queue consumed, so the
change is proportional to $x/D$: linear, with a slope inversely proportional to depth.
\iqlookfor{the four terms and the depth argument.}
\end{solution}

\begin{solution}{iq:rs:order-book-features:3}
An estimate of the fair price better than the mid, using the imbalance (and spread): the expected future mid given the book. Estimate by
bucketing the imbalance (and spread) and averaging subsequent mid changes, one step or iterated to convergence, on past data; validate out of
sample against the mid and the weighted mid.
\iqlookfor{conditioning on the state, estimation out of sample, the comparison.}
\end{solution}

\begin{solution}{iq:rs:order-book-features:4}
One reference implementation and a shared fixture: replay the same recorded messages through each implementation and require identical
outputs (to the last bit where possible, to a tight tolerance otherwise) in continuous integration, with the same operation order in the
arithmetic.
\iqlookfor{a replay test on a shared fixture, not code review.}
\end{solution}

\begin{solution}{iq:rs:order-book-features:5}
Size shown and withdrawn without intent to trade moves imbalance and depth features (spoofing, layering); flow features count those additions
as demand. Robustness: weight quantities by how long they rest or by the probability they trade, discount levels beyond the touch, use
executed flow as well as quoted size, and monitor for sequences of large adds and cancels.
\iqlookfor{the mechanism and at least two defences.}
\end{solution}

\begin{solution}{iq:rs:order-book-features:6}
P(up) is the probability that the ask queue empties first, increasing in the bid size and decreasing in the ask size; with equal rates it is
one half at equal sizes. It leaves out the efficient price (informed flow and stale-quote cancellations that make depletions dependent on
it), queue-reactive rates, and new levels inside the spread; it is therefore overconfident at extreme imbalances.
\iqlookfor{the race, its monotonicity, and what makes it overconfident.}
\end{solution}
